\section{Late exposure and the final matching}
\label{sec:late-dynamics-and-its-transfer-problem}

Write \(S_{\rm fin}\) for the cell-state configuration after the \(T_s\)
resampling rounds of Section~\ref{sec:initial-list-gates-and-resampling-estimates}.
These post-round states have not yet been conditioned on the terminal
avoidance requirements introduced below. The remaining task is to assign
the late odd roles injectively while preserving the common-neighbor lists
at even roles, and then to choose distinct labels from those lists.

\paragraph{Which law is being used.}
The argument uses the following processes and comparison laws.
\begin{enumerate}
\item \emph{Initial states.} The initial process draws permutation pools
and runs the finite resampling procedure to produce \(S_{\rm fin}\).
Terminal avoidance later conditions this same process. For nonnegative
tests on the specified local regions and deterministic horizons,
expectations are bounded using independent fresh cells at the same fixed
pools while the required pool conditions are retained, and then using fresh
cells on iid uniform slots.
\item \emph{Reference late transitions.} Given one entering history, each
late class uses independent row outputs consisting of masks, sketches, and
labels. The mask profiles are fixed independently of sampled histories.
The local gates are evaluated on actual data and the specified reference
kernels, and are retained as indicators in comparisons.
\item \emph{Actual late assignments.} Given a history satisfying the
incoming bounds and column-load condition, a conditional injective sampler
assigns the current class while avoiding its bad events and future alarms.
For the specified row tests, its upper comparison is to that class's
product reference law, and its normalizer applies only to the current draw
at the fixed entering history.
\end{enumerate}
For choosing the mask profiles, the annealed baseline averages over iid
slots, the full initial process including resampling and fallbacks, and
unconstrained product reference late transitions. Thus the iid input there
refers to the slots; the cell states still come from the full initial
process. The profiles are fixed under this baseline before terminal
conditioning.

\subsection{Reference transitions and local requirements}
\label{sec:reference-late-transitions}

\paragraph{Initial and current lists.}
Initial validity at an even \(w\) requires a valid internal prior
\(\sigma_w\) as above (a probability on initial internal hits, with its
indicated cap and single-corner cleaned support), \(S_w\) false including
the masked-list bound, and the indicated starting palette mass
\(\ge1/(4\chi_{i(w)})\). These hold everywhere after the desired initial
successes. They are gate predicates evaluated on the data. The reference
computations below are defined from direct data, with local fallbacks only
for invalid or zero computations.

Let \(\mathcal P_w^X\) be the host palette assigned to \(w\). On initial validity define
\[
 \tau_w^r(x)=\frac{U_w^0(x)\mathbf1_{\{x\in\mathcal P_w^X\}}}
 {\sum_z U_w^0(z)\mathbf1_{\{z\in\mathcal P_w^X\}}},
 \qquad \mathcal D_w^0=\operatorname{supp}\tau_w^r.
\]
Number the late classes \(r,\ldots,1\), so that \(j\) classes remain at
stage \(j\). The current probability \(\tau_w^j\) is \(\tau_w^r\)
conditioned to hit every previously assigned late neighbor of \(w\).
It reads only the initial data in \(w\)'s own slice, the external early
labels at \(w\), and its previous late neighbor labels. In particular,
it does not read the masks or sketches of previous late draws.

All reference kernels are defined on invalid histories as well. Fix a
label on each relevant host side in advance; an invalid initial list or
a zero normalizer is replaced by the point mass at that label. Use the
same lookup on every subsequent occurrence of the same direct data.
This choice affects no gated calculation below. More generally any fixed
probability fallback with these same inputs gives the comparisons used
here. Let \(p_w(j)\) count the neighbors of \(w\) in the remaining
classes. For a sufficiently small fixed \(\beta>0\), put
\(i=i(w)\) and define
\[
 e_{w,j}=C_n^{-\beta}\exp(-\beta s_i)2^{-\beta j}.
\]
The inequalities \(C_n\le n\), \(s_i=O(\log n)\) and
\(r=O(\log n)\) give \(\min e_{w,j}\ge n^{-.02}\) when \(\beta\)
is chosen after the low-mode constants. Also
\[
 \sum_{j=1}^r e_{w,j}
 \le K_\beta C_n^{-\beta}e^{-\beta s_i}=o(1).
\]
This remains \(o(1)\) after multiplication by \(\max(1,h_i)\): in
cluster mode \(s_i=(a/10^6)h_i\), and \(h_i e^{-\beta s_i}\) is
bounded even if the scale does not grow.

\paragraph{The product reference transition.}
Fix the entering history for class \(j\). The reference transition is a
product of independent row draws under that conditional law. At a late odd
role \(b\) of this class, draw the row as follows.
\begin{enumerate}
\item A mask profile is a probability distribution, fixed in advance,
over subsets of the current late pool of size at least half that pool.
Draw a subset from this profile and let \(\mu_{\rm mask}\) be its uniform law. For every even
\(w\sim b\), draw an iid sketch of \(\tau_w^j\) of length
\(\lceil n^{.25}\rceil\). These side draws are independent given history.
\item Test a candidate \(y\) by requiring empirical sketch hit fraction
\(\ge .5-e_{w,j}\) for each neighbor \(w\). Sample \(y_b\) by the masked
law conditioned on passing all tests if their retained mass is at least
\(\exp(-g_0n)\), \(g_0=\alpha/100\); otherwise just use the masked law.
\end{enumerate}
Write \(K_b^j\) for this full row kernel and
\(K_b^j(\,\cdot\mid\text{side data})\) for its conditional label law.
The mask and sketches belong to the side data in this output. A deletion kernel, given
these same side values, uses the same threshold rule after omitting the
indicated tests; it does not redraw the side data. No success checks are
inputs to these draws themselves. Masks with prescribed independent
profiles (specified for balance later) do not adapt to sampled histories.

\paragraph{Desired local conclusions.}
We want true hit mass for \(y_b\) under each \(\tau_w^j,\ w\sim b\), at least \(.5-3e_{w,j}\). We also want, given actual mask/sketch values, pointwise label-kernel upper comparison to each specified deletion kernel at factor \(\exp(1000 l e_{w,j})\), where one deletes any one external neighbor test (\(w\) that neighbor, \(l=1\)), or all internal neighbor tests at \(b\) (\(w\) any such internal neighbor, \(l=h_{i(b)}\)). Internal/external refers to the edge; an internal neighbor lies in the same patch.

\paragraph{The prerequisite gate.}
Let \(G_b\) be the following predicate of the entering history: initial
validity holds at every even site at Hamming distance at most \(6r\)
from \(b\), and the true-hit and deletion conclusions hold at every
earlier late role in that radius. We retain \(\mathbf1_{G_b}\) in each
bad event below. These history gates check actual data:
they evaluate the specified \textbf{reference} conditional kernels at
those data, without recursively verifying prior gate statuses. Thus raw
late-event simulation from \(S_{\rm fin}\) needs only direct data within
\(O(r)\) spatial radius, and similarly for forward influences (initial
direct data may of course require cell consultations there). On the
prerequisite gate, summable losses and at most one late hit per class at
\(w\sim b\) give a valid probability with
\begin{equation}
 \max\tau_w^j\le 8K C_n^{-1}\exp(-199s_i)2^{-p_w(j)}                   \label{eq:source-25}
\end{equation}
for \(K\) as in the initial cap. Indeed a previously imposed hit costs
at most \((.5-3e_{w,s})^{-1}\). Their product is at most
\(2^{p_w(r)-p_w(j)}\exp(O(\sum_s e_{w,s}))\); the last factor is at
most two eventually.

The three requirements are as follows.
\paragraph{Requirement 1 (sketch conflicts).}
For a sketch \((X_1,\ldots,X_m)\) of \(\tau_w^j\), where
\(m=\lceil n^{.25}\rceil\), require
\[
 m^{-2}\sum_{a,c\le m}\mathbf1_{\{X_a,X_c\text{ conflicting}\}}
 \le2e_{w,j}^4.
\]
The ordered sum includes its diagonal.
\paragraph{Requirement 2 (prefix moments).}
Fix an order with the external tests singly, followed by the internal
tests as one batch; omit an empty internal batch. Include also the orders
obtained by moving any specified external test to the end. For a batch
in one of these orders let \(A_q\) be the intersection of the preceding
sketch tests, excluding the current batch. Whenever
\(\mu_{\rm mask}(A_q)\ge e^{-g_0n}\), set
\(U=\mu_{\rm mask}(\,\cdot\mid A_q)\). Thus \(U\) is a law on
late host labels. For each even role \(v\) in the current batch
require
\begin{equation}
 \int I_x\,dU=.5\pm10 b_*,\qquad \int I_x I_z\,dU=.25\pm10 b_*         \label{eq:source-26}
\end{equation}
for all \(x\in\mathcal D_v^0\) and all nonconflicting pairs there.
\paragraph{Requirement 3 (true hits).}
Require \(d_G(\tau_w^j,y_b)\ge .5-3e_{w,j}\) for every \(w\sim b\).

The sketch-conflict and prefix-moment bad events are their respective
failures intersected with \(G_b\). The true-hit bad event is
\(G_b\cap\{\text{Requirements 1 and 2 hold}\}\) intersected with a
failure of Requirement 3. These are events in the reference process,
not new conditionings of its kernels.

The next lemma also records the single-label and pair-hit kernel bound used
later in the endpoint comparison.

\begin{lemma}[Consequences of the local transition requirements]\label{lem:late-local-transition}
Fix an entering history satisfying the prerequisite gate, and use the product reference transition at \(b\). There is a fixed \(c>0\) such that, uniformly over the fixed mask profiles,
\[
\begin{aligned}
 \Pr(\text{Requirement 1 fails})&\le e^{-n^c},\\
 \Pr(\text{Requirements 1 and 2 hold but a true-hit conclusion fails})
       &\le e^{-n^c}.
\end{aligned}
\]
If the first two requirements hold, every tested prefix is broad, the specified deletion comparisons hold, and the full conditional label kernel satisfies \eqref{eq:source-27}. The true-hit conclusion is enforced by the third requirement.
\end{lemma}

\Needspace{5\baselineskip}
\begin{proof}\leavevmode\par
\paragraph{Sketch concentration.}
The conditional expected conflict fraction is at most
\[
 m^{-1}+\Delta_i\|\tau_w^j\|_\infty
 \le m^{-1}+8K^2C_n^{-1}e^{-198s_i}2^{-p_w(j)}
 \le e_{w,j}^4
\]
eventually. The last inequality uses \(p_w(j)\ge\lfloor j/2\rfloor\)
and small \(\beta\); the conflict degree is used only on the active
single-corner support. Write \(m=\lceil n^{.25}\rceil\) for the sketch length and \(e=e_{w,j}\) for a fixed target. Changing one sample changes the empirical conflicting-pair fraction by at most \(2/m\), so its excess above \(2e^4\) costs \(\exp(-\Omega(m e^8))\) by bounded differences. Since \(e\ge n^{-.02}\), this is \(\exp(-\Omega(n^{.09}))\). A union over at most \(n\) targets gives the first asserted bound for some fixed \(c>0\).
\paragraph{Broad prefixes and deletion kernels.}
Suppose Requirements 1 and 2 hold. For a target sketch, its empirical
hit fraction \(H(y)=m^{-1}\sum_a I_{X_a}(y)\), with \(y\sim U\),
has mean \(.5+O(b_*)\). In its second moment, nonconflicting pairs
contribute \(.25+O(b_*)\) by \eqref{eq:source-26}, and conflicting
pairs have total fraction at most \(2e_{w,j}^4\). Hence
\[
 \operatorname{Var}_U H=O(b_*+e_{w,j}^4),\qquad
 U\{H<.5-e_{w,j}\}
 \le O\big((b_*+e_{w,j}^4)/e_{w,j}^2\big)=O(e_{w,j}^2).
\]
Here \(b_*\ll e_{w,j}^4\). Multiply the retention factors for external
tests, and use a union bound for the internal batch. The total retained
mass is \(\exp(-o(n))\), and the internal batch loses \(o(1)\).
Starting with the masked law, this proves inductively that every tested
prefix exceeds \(e^{-g_0n}\). Thus neither the full label kernel nor
the relevant deletion kernel uses its fallback.

Place the deleted external test last, or use the base order with the
internal batch last. The density of the full kernel relative to the
leave-out law is the reciprocal retained mass of that last test or
batch, at most \(\exp(1000l e_{w,j})\). Applying
\eqref{eq:source-26} under this leave-out law now bounds a single hit
or a nonconflicting pair hit from \(\mathcal D_w^0\) by
\begin{equation}
 2^{-u'}\exp(O(l e_{w,j})),\qquad u'=1,2                              \label{eq:source-27}
\end{equation}
respectively, where \(l\) is the size of the deleted external test or
internal batch.

\paragraph{True hits.}
On the first two requirements, labels meet the sketch tests, so compare to the corresponding deletion laws while retaining those requirements. The comparison label for \(w\) no longer depends on that neighbor's sketch. Its pointwise multiplier \(\exp(1000l e_{w,j})\) is bounded by the uniform error-sum estimate, including when \(l=h_i\). We may now drop the success restrictions in the nonnegative upper integral. For a comparison label whose true hit mass is below \(.5-3e_{w,j}\), the independent target sketch passes the threshold \(.5-e_{w,j}\) with probability at most \(\exp(-\Omega(m e_{w,j}^2))\le\exp(-\Omega(n^{.21}))\), by the bounded-summand estimate. The union over targets proves the second asserted bound, decreasing \(c\) if necessary. This bounds the intersection with Requirements 1 and 2; it does not condition the reference law on their success.
\end{proof}

Avoiding all these gated bads after globally good initial states suffices by induction (our stage assignments with joint upper comparison will indeed use only values in reference support given history). The failure still needing a transfer through fresh initial data is Requirement~2, namely \eqref{eq:source-26}: a broad \(U\) is history-adaptive and the supports have used early data.

\subsection{Transfer for a late prefix}
\label{sec:transfer-for-a-late-prefix}

We call the bad events just defined the late events \(F\). For
Requirement 2 use one event for each fixed role, target and prefix,
bundling all its single-label and pair tests. Fix such an event at
\(v\sim b\), with \(i=i(v)\). Let \(\mathcal I_i\) be the internal
coordinate set of patch \(i\), and \(t\) the word of \(v\) on its
complement. Throughout this proof outer words are relative to
\(\mathcal I_i\), even at sites in other patches. Write \(A_k\) for
the coset modulo \(\mathcal L\) of coordinate \(k\)'s ID, \(C_w\)
for the syndrome coset of a cube vertex \(w\), and \(w^k\) for a
coordinate flip.

The critical cells are those containing \(v^k\) for bulk coordinates
\(k\) with \(A_k\ne C_v\). These are distinct cells of patch \(i\),
and the indicated roles are external early neighbors of \(v\), in
slices with outer words \(t^k\). There are \(d=n-O(\log n)\) of
them. One such cell may be omitted; if so, put it among the fixed
noncritical inputs and impose no witness hit from it in the estimate.
The experiment for the lemma is:
\begin{itemize}
\item fix arbitrary initial states in all other needed cells;
\item draw each critical whole-cell state independently, using an
unpinned iid slot pool conditioned on individual typicality and then a
fresh cell sample on that pool;
\item run the product reference late transitions from these states.
\end{itemize}
On the gated prefix failure there is a label \(x\in\mathcal D_v^0\),
or a nonconflicting pair \(x,z\in\mathcal D_v^0\), violating
\eqref{eq:source-26}. Every retained critical label hits every label
in this witness, because \(\mathcal D_v^0\) satisfies all early hits
at \(v\). This necessary condition is what permits us to sum over
possible witnesses below.

\begin{lemma}[Late-prefix transfer]\label{lem:late-prefix-transfer}
In the preceding experiment, the fixed prefix-moment bad event \(F\) has probability at most \(e^{-n^{c_1}}\), for a fixed \(c_1>0\). The estimate is uniform over the fixed states of all noncritical cells and the independent mask profiles, and still holds when one critical cell is omitted as allowed above.
\end{lemma}

\begin{proof}\leavevmode\par
The proof first compares reference path factors with deletion kernels on
the required path event, and then integrates the erased sketches. Only
afterward does it fix independent late-draw seeds and construct a
deterministic block simulation. A survival tilt and a stopped likelihood
estimate then bound the rare prefix deviation.

\paragraph{Critical labels.}
Let \(y_k=y_{v^k}\), and denote its marginal law by \(q_k\). Integrating
the fresh singleton comparison over an unpinned iid pool conditioned on
individual typicality gives
\(q_k\le(1+O(n^{-3}))\pi_i\), and therefore total-variation distance
\(O(n^{-3})\). Direct spatial consultations for this event lie within
\(O(r)\) of \(b\). Cell spacing implies that they reach each critical
cell only in its indicated slice. This does not assert independence of
that slice from the other slices of its cell.

\subsubsection{Erasing part of a history}
\label{sec:erasing-part-of-a-history}

\paragraph{The predecessor set.}
Start with the mask and prefix sketches at \(b\). Whenever a needed
sketch at an even site \(w\) uses a previous late neighbor's label,
include that neighbor's full reference transition and all sketches
needed by it. Continue backwards through the class order, generating
each late role once. Call these the direct roles, and call the even
sites queried by their sketches the direct even sites. This finite set
is closed under label predecessors, so its marginal law is exactly the
product of the corresponding conditional reference kernels. Gate-only
consultations outside it may be discarded when taking an upper bound.

There are at most \(r\) backward class steps, so every direct role is
within distance \(2r\) of \(b\). Let \(P(w)\) record, over the binary
field, the parity of the number of active coordinates in each ID coset.
A two-edge passage between late roles flips two IDs in the same coset,
because both late syndromes belong to \(\mathcal L\). Consequently
\(P(b')=P(b)\) for every direct late role \(b'\), and
\[
 P(w)=P(b)+\boldsymbol e_{C_w},\qquad
 P(w)+P(v)=\boldsymbol e_{C_w}+\boldsymbol e_{C_v}
\]
for every direct even site. The prerequisite gate supplies initial
validity, positive current lists and earlier deletion comparisons on
this whole set, which lies within its stated radius.

\paragraph{Erase the exposed outer word.}
Let \(H_i=\{A_a:a\in\mathcal I_i\}\); these cosets are distinct.
If \(D=\{q\notin\mathcal I_i:w_q\ne v_q\}\), the parity identity is
\begin{equation}
 \sum_{\substack{a\in\mathcal I_i\\w_a\ne v_a}}\boldsymbol e_{A_a}
 =\boldsymbol e_{C_w}+\boldsymbol e_{C_v}
       +\sum_{q\in D}\boldsymbol e_{A_q}.
 \label{eq:late-coset-parity}
\end{equation}
Its right side must be supported on \(H_i\), and then uniquely
determines the internal bits of \(w\).

Let \(V_0\) consist of direct even sites with outer word \(t\).
If \(C_v\notin H_i\), setting \(D=\varnothing\) in
\eqref{eq:late-coset-parity} forces \(C_w=C_v\) and \(w=v\).
If \(C_v\in H_i\), let \(k_0\) be its unique internal coordinate.
The same identity gives \(w=v^{k_0a}\) for some internal \(a\), so
\(V_0\) is contained in the internal neighbors of \(b_0=v^{k_0}\).
A direct late role with outer word different from \(t\) can meet
\(V_0\) only through the unique external flip returning its word to
\(t\). A direct late role with outer word \(t\) must be \(b_0\),
by \(P(b')=P(b)\), and all its internal tests lie in \(V_0\).

At each earlier direct role meeting \(V_0\), replace its label kernel
by the deletion kernel omitting that external test or internal batch.
Keep its side-draw factors unchanged. On the original gated failure,
the earlier deletion comparisons bound the product of original path
factors by the product of modified factors times
\[
 \exp\left(1000\max(1,h_i)\sum_s e_{v,s}\right)=O(1).
\]
Indeed at most \(\max(1,h_i)\) tests are deleted in each class. The
comparison is made at path values, before fixing any randomization seeds.
At the current role \(b\), the tested prefix already excludes the
single target \(v\) in the external case, or the entire internal
batch in the other case. It therefore contains no sketch from \(V_0\).

\paragraph{The event retained after deletion.}
For each remaining direct sketch at \(w\), form the set \(L_w\) of
envelope labels at \(w\) hitting all its unaffected external early
labels. Ignore the own prior, all critical external early inputs, and
all late hits in this definition. Thus \(L_w\) uses only fixed
noncritical data. We retain the following weaker event
\(\mathcal E_{\rm keep}\): these sets have size at most
\(\exp((\log n)^8)\) at every affected sketch site; each remaining
sketch there lies in \(L_w\); the tested prefix is broad; and there
is an envelope witness of the required kind which survives every
critical hit and violates \eqref{eq:source-26} for that prefix.
The one-block calculation below shows that the sets \(L_w\) omit at
most \((\log n)^4\) external early inputs. Thus the initial
masked-list requirement and the supported current lists on the
original gate imply these retained conditions.

The original failure is contained in \(\mathcal E_{\rm keep}\),
after the pointwise comparison has been applied. This weaker event
reads only initial data, generated labels and remaining sketches. The
erased sketches no longer enter any label kernel, and earlier side
outputs are never inputs to later current lists. They enter neither
the retained event nor a remaining transition, so their conditional
probability factors integrate to one. In particular no deletion test
on an erased sketch is being retained implicitly.

\paragraph{One-block locality.}
Partition the critical coordinates by their ID cosets, merging all
cosets in \(H_i\) into one block. Each block has at most
\(m_*:=\max(1,h_i)r\) coordinates. We show that a remaining direct
sketch consults at most one such block.

If its own slice is critical, its outer word is \(t^k\). An external
flip produces either \(t\), whose sketches have been erased, or a
word \(t^{kj}\) differing in two external coordinates; it cannot
reach another critical slice. Internal flips stay inside the same
patch and belong to the own slice computation. Hence this case reads
only the block containing \(k\).

Otherwise a critical external observation at \(w^j\), in slice
\(t^k\), requires outer word \(t^{kj}\) at \(w\), with
\(j\notin\mathcal I_i\) and \(j\ne k\). An internal flip would
stay in that same patch, whereas \(j=k\) would put \(w\) in the
erased word \(t\). There can be only one further critical observation:
flipping \(k\) reaches \(t^j\). If both observations are critical,
then \(A_k,A_j\ne C_v\); their being early gives
\(A_k,A_j\ne C_w\). If \(A_k\ne A_j\), neither coordinate can
cancel on the right of \eqref{eq:late-coset-parity}. Both therefore
belong to \(H_i\). They lie in the merged block; if their cosets
are equal, they already lie in one block.

This calculation also bounds the number of affected sites. For fixed
critical \(k\), an external observation through \(j\) has
\(A_j\ne C_w\). Equation~\eqref{eq:late-coset-parity} then forces
\(A_j\in H_i\cup\{A_k,C_v\}\), giving at most \((h_i+2)r\)
choices of \(j\). For each resulting outer word, at most \(h_i+3\)
values of \(C_w\) are possible, each determining all internal bits.
A block of size \(m\) therefore affects at most
\[
 m\,[1+(h_i+2)r](h_i+3)
\]
sites, with at most \(r\) sketch calls at each. These counts are
polylogarithmic in \(n\). They include cross-patch observations:
\(\mathcal I_i\) lies beyond every allocation prefix, so the outer
word fixes the patch, and the parity identity remains global.

\subsubsection{A bounded-response simulation}
\label{sec:a-bounded-response-simulation}

\paragraph{A total protocol after fixing seeds.}
Now fix independent seeds for all late draws in the modified process,
independently of the initial states. Order the finite predecessor
computation by class and then by a fixed role order. A block input is
the collection of whole critical cell states in that block, including
their pools. At an affected sketch call, send that block the preceding
labels and fixed data needed for the current list. The block computes
the sketch using the fixed seed. It replies with the sketch's
\(\lceil n^{.25}\rceil\) indices in the fixed list \(L_w\), if
\(|L_w|\le\exp((\log n)^8)\) and every sampled label belongs to it.
Otherwise it replies with a single abort symbol. This is a total rule:
the current-list fallbacks were fixed when defining the reference law.

On receiving an abort, halt and output the current role \(b\)'s
masked law. Its mask is determined by its independent fixed seed. In
the absence of an abort, use the replies to evaluate the remaining
modified label transitions. At the end output the prefix-conditioned
law if its mass is at least \(e^{-g_0n}\), and the masked law otherwise.
Unaffected computations use only fixed noncritical inputs and preceding
replies. Every request is therefore determined without reading any
other block. On \(\mathcal E_{\rm keep}\) there is no abort, and
induction in the chosen order shows that every reply and generated
label agrees with the modified reference process. Its final output is
the tested prefix law.

The list-size assertion used here follows from the preceding geometry:
\(m_*=O((\log n)^{1.1})\le(\log n)^4\) eventually, so omitting
a block omits no more early inputs than the initial masked-list bound
allows. The output \(U\) of the protocol is always broad, with
\[
 \operatorname{width}(U)
 \le\log(2N/M_{\rm late})+g_0n
 =O(\log r)+g_0n<\alpha n/2.
\]

\paragraph{Information returned by each block.}
Only \(n\,\mathrm{polylog}(n)\) response steps are needed. This
counts calls involving critical blocks, not every node of the possibly
larger predecessor computation. Each call returns a sketch with at most
\(n^{.25+o(1)}\) indices, and each index uses at most \((\log n)^8\)
logarithmic range. The number of calls to any one block is
polylogarithmic by the site count above. Including the abort symbol,
the entire reply sequence of one block has log-range at most
\(n^{.4}\). No size bound on requests is needed.

Fix this protocol, including its fallbacks and response alphabets,
before choosing witness labels. The event \(\mathcal E_{\rm keep}\)
is now contained in the event that some envelope label, or
nonconflicting envelope pair, survives all critical hits and violates
\eqref{eq:source-26} for the protocol's broad output. We bound this
larger event. Its protocol is independent of the witness being tested.

\paragraph{Raw states and survival conditioning.}
Integrate a witness \(\boldsymbol x\) against the uniform law on
\(X_i^{a'}\), where \(a'=1\) or \(2\). Let \(\mathcal A\)
require envelope membership and, for a pair, nonconflict. Write
\(\mathbb P_0\) for the raw independent block law. For a block \(H\),
let \(A_H(\boldsymbol x)\) be the event that all its critical labels
hit every witness label, and put
\[
 c_k(\boldsymbol x)=\Pr_{q_k}(y_k\text{ hits }\boldsymbol x),\qquad
 c_H(\boldsymbol x)=\mathbb P_0(A_H)=\prod_{k\in H}c_k(\boldsymbol x).
\]
On \(\mathcal A\), each \(c_k\ge .15\). For pairs this follows from
\[
 \int I_x I_z\,d\pi_i
 =\tfrac14(1+m_{\pi_i}(x)+m_{\pi_i}(z)+K_{\pi_i}(x,z)),
\]
the cleaned degree bounds, \(|K_{\pi_i}(x,z)|\le\xi\), and the
\(O(n^{-3})\) perturbation from \(\pi_i\) to \(q_k\). The
single-label case follows directly from the degree bounds.

For fixed allowed \(\boldsymbol x\), let \(\mathbb P_{\boldsymbol x}\)
condition each whole block independently on \(A_H(\boldsymbol x)\).
Equivalently it conditions each critical whole-cell state, including
its pool, on the required hits of its critical label. We first prove
that under this law the probability of witness deviation into the
protocol's output, averaged over \(\boldsymbol x\) with
\(\mathbf1_{\mathcal A}\), is \(\exp(-n^{\Omega(1)})\).

\paragraph{Transcript cylinders.}
For a transcript prefix \(t\), let \(C_H(t)\) be the set of block
inputs giving its prescribed replies to the prescribed requests.
Since requests depend only on preceding replies and fixed data,
\[
 \mathbf1_{\{T=t\}}=\prod_H\mathbf1_{C_H(t)}.
\]
Conditional on this prefix, raw blocks are therefore independent,
each restricted to its own cylinder. Fixing all other blocks, one
block's reply sequence determines the full transcript. There are at
most \(e^{n^{.4}}\) possible such sequences. At any specified update,
the raw probability that the selected block's cylinder has mass below
\(e^{-\alpha n/8}\) is consequently at most
\(e^{n^{.4}-\alpha n/8}\), with a further polynomial factor harmless
if the block is selected adaptively. This is an integrated raw bound,
not a bound conditional on every transcript.

Off this exception, if the block has \(m\) critical labels their joint
conditional law satisfies
\[
 Q\le e^{\alpha n/8+o(n)}\pi_i^{\otimes m}.
\]
The unconditioned tuple has a product upper bound because its labels
belong to distinct independent cells. Independence of those labels is
not asserted after restriction to a cylinder.

\paragraph{Survival under a cylinder.}
For any such \(Q\), let
\(q_j=dQ_{\le j}/d\pi_i^{\otimes j}\). The \(Q\)-mass of value
prefixes with \(q_j<e^{-\alpha n/8}\) is at most
\(e^{-\alpha n/8}\). On the other prefixes the next-coordinate
density relative to \(\pi_i\) is at most
\(e^{\alpha n/4+o(n)}\), leaving its width below \(\alpha n/2\).
For any second-side law with that width, deep discrepancy implies
that a uniform \(x\in X_i\) has degree \(.5\pm2b_*\), except
on mass \(\exp(-\Omega(n^{x_*}))\). Indeed a larger exception of
one sign could be normalized to a first-side law within the other
budget of \eqref{eq:source-9}. For two witness labels, first impose
the typical hit of \(x\), then apply the same estimate to the
normalized restriction for \(z\). This costs only bounded width.

Thus each next coordinate survives with probability
\(2^{-a'}+O(b_*)\), apart from a tiny set of witness tuples for each
fixed value prefix. Fubini bounds the mean, over witness tuples, of
the \(Q\)-mass of exceptional value prefixes. Markov and a union
over the \(m\) coordinates make this mass
\(\exp(-n^{\Omega(1)})\) outside witness mass of that same form.
The small-density prefixes have already been bounded in absolute
\(Q\)-mass. For a witness outside the resulting exceptional set, let
\(A_{\le j}\) mean that the first \(j\) coordinates hit every witness
label, with \(A_{\le0}\) the whole space. At each step the bad value
prefixes have total \(Q\)-mass at most
\(\varepsilon_n=\exp(-n^{\Omega(1)})\). Integrate the conditional hit
bound with the indicator of \(A_{\le j-1}\), bounding its contribution
on those bad prefixes by their total mass. This gives
\[
 Q(A_{\le j})
 =\bigl(2^{-a'}+O(b_*)\bigr)Q(A_{\le j-1})+O(\varepsilon_n).
\]
Since \(m\le m_*\) is polylogarithmic, iteration makes the additive
errors negligible relative to \(2^{-2m}\), and gives
\[
 Q(A_H(\boldsymbol x))
 =2^{-a'm}(1+O(mb_*))
\]
outside an exceptional uniform witness mass \(\exp(-n^{\Omega(1)})\).
Apply this estimate both to a transcript-conditioned block and to its
raw law before any transcript. The latter application controls the
denominator \(c_H\); the cutoff alone supplied only its positive
lower bound \(.15^m\). The transcript-conditioned estimates are
averaged over the raw transcript distribution. There is no union over
all transcripts.

\paragraph{The likelihood process.}
Include the fixed witness in the transcript filtration
\(\mathcal F_t\), and define
\[
 r_H^t=\frac{\mathbb P_0(A_H(\boldsymbol x)\mid C_H(T_t))}{c_H(\boldsymbol x)},
 \qquad R_t=\prod_Hr_H^t
 =\frac{d\mathbb P_{\boldsymbol x}(T_t)}{d\mathbb P_0(T_t)}.
\]
For fixed allowed witness, \(R_t\) is a nonnegative martingale
under \(\mathbb P_0\), starting at one. At an update only the responding block changes.
The preceding estimates give \(r_H^t=1+O(mb_*)\), except on an event
whose raw probability averaged with \(\mathbf1_{\mathcal A}\)
is at most \(e^{-n^{c_2}}\), for a fixed \(c_2>0\). This exceptional
event includes a small cylinder, a failed conditional survival
estimate, or a witness with a failed no-transcript denominator estimate.
The last set has small integrated mass even though every ratio starts
at one on it.

Stop at the first exception to these factor bounds or the first
\(R_t>B_*:=e^{n^{c_3}}\). Choose \(c_3>0\) smaller than \(c_2\)
and than the positive exponent in the broad-law deviation estimate
below. Before stopping, the next multiplier
\(Z_t=r_H^t/r_H^{t-1}\) has conditional mean one and is
\(1+O(m_*b_*)\) off the next exception. On that exception,
\(c_H\ge .15^m\) and the previous good factor give the crude
bound \(Z_t\le e^{O(m_*)}\). Hence
\[
 Z_t^2\le1+2(Z_t-1)+O(m_*^2b_*^2)
                   +e^{O(m_*)}\mathbf1_{\rm next\ exception}.
\]
Writing \(\tau\) for the stopping time and averaging over both raw
states and uniform witnesses gives
\begin{align*}
 \mathbb E[\mathbf1_{\mathcal A}R_{t\wedge\tau}^2]
 &\le (1+O(m_*^2b_*^2))
       \mathbb E[\mathbf1_{\mathcal A}R_{(t-1)\wedge\tau}^2]\\
 &\hspace{1cm}+B_*^2e^{O(m_*)-n^{c_2}}.
\end{align*}
The linear term vanishes by the martingale property. The barrier is
what permits the merely integrated exception estimate in the last
term. There are at most \(n\,\mathrm{polylog}(n)\) updates, and
\(n\,\mathrm{polylog}(n)b_*^2=o(1)\); thus the stopped second
moment is \(O(1)\).

Under the tilted law, a barrier stop has integrated probability at
most \(B_*^{-1}\mathbb E[\mathbf1_{\mathcal A}R_\tau^2]=O(B_*^{-1})\).
Exception stops cost \(\exp(-n^{\Omega(1)})\), using the pre-stop
barrier, the maximum next multiplier and the raw exception bound.
Off stops, tilted deviation probability is at most \(B_*\) times
raw deviation probability. For every fixed transcript the output
\(U\) is broad and independent of the witness; deep discrepancy,
applied sequentially for a pair, bounds its exceptional uniform
witness mass by \(\exp(-n^{\Omega(1)})\), with the slack in
\eqref{eq:source-26}. The choice of \(c_3\) leaves the same type of
bound after multiplication by \(B_*\). This proves the required
integrated tilted estimate.

\paragraph{Returning to surviving witnesses.}
Let \(p_{\rm tilt}(\boldsymbol x)\) be the tilted deviation
probability. Undoing the survival conditioning supplies the factor
\(\prod_k c_k(\boldsymbol x)\). The integrated singleton comparison,
the low-mode degree bounds and the second-moment estimate after
\eqref{eq:source-24} give
\[
 \mathbb E_{\boldsymbol x}\!\left[
 \mathbf1_{\mathcal A}(2^{a'd}\prod_k c_k(\boldsymbol x))^2\right]
 \le e^{O(\log n)}.
\]
For a singleton this uses the degree bounds; for pairs it uses the
stated pair second moment. The relative \(1+O(n^{-3})\) perturbation
of each critical marginal contributes a bounded factor.
Cauchy--Schwarz, using \(p_{\rm tilt}^2\le p_{\rm tilt}\), bounds
the expected number of surviving deviating witnesses by
\[
 M_i^{a'}2^{-a'd}
 \left(\mathbb E\mathbf1_{\mathcal A}
            (2^{a'd}\prod_k c_k)^2\right)^{1/2}
 \left(\mathbb E\mathbf1_{\mathcal A}p_{\rm tilt}\right)^{1/2}
 =\exp(-n^{\Omega(1)}).
\]
Indeed \(M_i\le n2^n\) and \(d=n-O(\log n)\), so the first factor
is polynomial. The bounds are uniform in the fixed seeds and
noncritical states. Averaging them and restoring the bounded path
comparison proves the lemma.
\end{proof}

\subsection{Terminal conditioning of initial sampling}
\label{sec:terminal-conditioning-of-initial-sampling}

Fix any family of independent mask profiles of the kind specified above.
Given \(S_{\rm fin}\), let \(p_F\) be the probability of each late bad \(F\)
under the reference late process for these profiles, with no late constraints
imposed. Include as terminal requirements:
\begin{itemize}
\item individual pool typicality and \eqref{eq:source-23} everywhere;
\item final \(S_v\) avoidance;
\item \(p_F\le e^{-n^\delta}\) for every \(F\), with sufficiently small fixed \(\delta>0\).
\end{itemize}

For each of the latter two kinds, gate the failure on the pool requirements throughout a deterministic region sufficient for its restricted resampling simulation. That is, start with all involved direct cell consultations and their incidence events (also the defining site as needed), then take simulation balls of radius \(2T_s+O(1)\) large enough by deterministic locality, including all required event scopes and pool inputs. On success of the first kind globally these gates are passed. Slot and cell-tape scope sizes per requirement, and requirement incidences per slot or cell tape, are at most
\(n^{O(A_c+4)T_s+O(r)}\). Indeed graph neighborhoods per resampling round enlarge by polynomial factors as above. Before expanding those simulations the raw \(F\) consultations/influences only use spatial radius \(O(r)\) (including paths to compute direct gate conditions), with polynomial cell/position size ratios and polynomially many tests per late role. Here the absolute constants fit by taking \(P\) sufficiently large in terms of \(A_c\) as above, using \(r=o(T_s)\).

Write \(\mathcal G_F\) for this local pool gate when considering the
terminal requirement for \(p_F\).

\begin{proposition}[Terminal avoidance and local comparison]\label{prop:terminal-avoidance}
The terminal requirements admit simultaneous conditioning on avoidance. Each bad requirement has probability at most \(n^{-PT_s/3}\) before conditioning, also after any one global slot pin. Under the resulting law, nonnegative tests on the stated unions of local regions have upper comparison with unconditioned initial sampling at a multiplicative cost tending to one.
\end{proposition}

\begin{proof}\leavevmode\par
\paragraph{Bounds with one prescribed slot.}
Pool failures and final \(S_v\) failures already have the required
\(n^{-PT_s/3}\) bound from Lemma~\ref{lem:uniform-pool-list-estimate}
and Proposition~\ref{prop:initial-resampling-comparison}, also after
one prescribed slot-to-bin mapping. For a late event define
\[
 p_F(S_{\rm fin})=\Pr_{\rm ref}(F\mid S_{\rm fin}).
\]
The sketch-conflict and true-hit events have the required tiny
expectation uniformly in the entering history by
Lemma~\ref{lem:late-local-transition}. The prefix-moment event
requires the following replay.

Use its critical cells, omitting the cell of the globally prescribed
slot if it is among them. Mark every resampling event touching a
critical cell and every neighbor of such an event in the cell-scope
graph. Let the target cells be the union of scopes of events touching
critical cells. Their number is polynomial, below
\(n^{10(A_c+10)}\). Every marked execution touches this target set,
so belongs to its backward execution closure. Let \(B_F\) be the
event that this closure has more than \(T_s\) executions. On the
required pool gate, its probability is at most \(n^{-PT_s/2}\).
Put \(\eta=e^{-n^\delta}\) and split before applying Markov:
\[
 \Pr(\mathcal G_F,\ p_F>\eta)
 \le \Pr(\mathcal G_F\cap B_F)
    +\eta^{-1}\mathbb E[
       \mathbf1_{\mathcal G_F}\mathbf1_{B_F^c}p_F].
\]
Thus the factor \(\eta^{-1}\) will apply only to the replay estimate
on \(B_F^c\).

\paragraph{Replay with a full neighbor buffer.}
On \(B_F^c\), record the complete list \(\omega\) of marked site/round
executions. At most \(T_s\) entries gives at most
\(\exp(O((A_c+4)T_s\log n))\) possible lists. For one fixed list,
force precisely those marked executions in their recorded rounds,
regardless of truth. At unmarked sites use the original priority rule,
including its readings of neighboring truths. If off-path forced and
unforced executions overlap, advance each involved cell only once
in that round. This defines a replay on every input and agrees with
the original process whenever its marked occurrence list is \(\omega\).

An event whose truth reads a critical cell is marked. Every decision
that could read that truth is marked as well, because the entire
neighbor buffer was included. Hence no unmarked decision reads a
critical truth. The final critical tape indices are prescribed by
\(\omega\); noncritical final outputs depend only on noncritical
inputs. For this fixed pattern, discard its consistency condition in
the nonnegative upper integral, retaining individual typicality of
critical pools. The consulted permutation slots compare to iid slots
at \(1+o(1)\) cost, with the global slot pin maintained outside the
critical cells. Conditional on the other inputs and on each critical
pool's individual typicality, its prescribed terminal entry is a
fresh cell state, independently across critical cells. Thus
Lemma~\ref{lem:late-prefix-transfer} gives
\[
 \mathbb E\big[\mathbf1_{\mathcal G_F}
                  \mathbf1_{\{\text{pattern }\omega\}}
                  p_F(S_{\rm fin})\big]
 \le (1+o(1))e^{-n^{c_1}}.
\]
The inequality uses the enlarged replay integral just described;
normalizing the individual typicalities costs only negligible factors.
Summing these pattern bounds in the truncated expectation above,
with \(\delta<c_1\), proves the terminal bad bound. Cell tapes may be
represented by independent randomizers before pools are drawn, using
fixed lookups for every possible local pool.

\paragraph{Leaves of the permutation experiment.}
Split each terminal bad event into leaves specifying the exact images
of its consulted permutation slots and a measurable event of its
consulted tapes. Declare two leaves neighbors if they share a domain
slot, an image bin, or a cell tape. Image collisions matter even when
their domain slots differ. For the negative-dependency precedent for
partial mappings in a uniform injection see Lu and Sz\'ekely
\cite[Theorem~1]{lu-szekely2007}. Here the following coupling also
handles the independent tape variables.

Force a positive-probability leaf by successively swapping each
prescribed image into its required slot, and condition its tape data
on the prescribed tape event. This yields the conditional law given
the leaf and preserves every occurring nonneighbor leaf: their domains,
images and tapes avoid those used in the forcing. Consequently
\[
 \Pr(\text{avoid nonneighbors}\mid L)
 \le\Pr(\text{avoid nonneighbors}),\qquad
 \Pr(L\mid\text{avoid nonneighbors})\le\Pr(L).
\]

Let \(D_n=n^{C(A_c+4)T_s+O(r)}\) bound both scopes and incidences,
for one fixed constant \(C\). At a fixed domain slot or tape the
sum of incident leaf probabilities is at most
\(D_n n^{-PT_s/3}\). For a fixed image bin \(D\) in patch \(i\),
its sum is at most
\[
 \sum_{E,s}\Pr(D_s=D,\ E)
 =B_i^{-1}\sum_{E,s}\Pr(E\mid D_s=D)
 \le D_n n^{-PT_s/3},
\]
where each event is charged through its consulted slots. There are
at most \(B_i\) such slots, each with incidence at most \(D_n\),
and the one-slot bound applies to every summand. Charges twice the
raw leaf probabilities satisfy the conditional local lemma when
\(P\) is sufficiently large compared with \(C(A_c+4)\), since a
leaf touches at most \(D_n\) slots, images and tapes.

The same argument compares tests supported on a polynomial union of
these local regions, or on \(n^{O(r)}\) direct consultations expanded
by their deterministic simulation horizons. Split a test event by its
consulted slot images. Each resulting leaf has the same nonneighbor
property; the sum of touching charges is \(o(1)\). Dropping those
constraints therefore costs a reciprocal avoidance product \(1+o(1)\).
Summing test leaves, and then integrating level sets, proves the
stated comparison for every nonnegative test on these deterministic
scopes. This includes the column moments and endpoint products below.
\end{proof}

\subsection{Staging and late injectivity}
\label{sec:staging-and-late-injectivity}

\begin{proposition}[Completion of the late assignment]\label{prop:late-assignment}
There are independent mask profiles fixed in advance and conditional injective samplers for the successive late classes such that the full run completes with probability tending to one under terminal-conditioned initial sampling. On a full run all true-hit and deletion conclusions hold. Each executed class has the localized upper comparison with its reference transition used below.
\end{proposition}

\Needspace{5\baselineskip}
\begin{proof}\leavevmode\par
\paragraph{Balance in the baseline.}
Use the full initial process, including finite resampling and fallbacks,
on iid uniform slots before terminal conditioning, followed by
unconstrained reference late transitions. This is the baseline for
choosing profiles; its cell states are not fresh independent states.
Fix profiles in processing order. For one row \(b\), the preceding
baseline history distribution is then fixed. Consider the convex hull
of its baseline label-marginal vectors obtained from the allowed masks.
Against any normalized nonnegative label prices, the cheapest
\(\lceil M_{\rm late}/2\rceil\) labels all have price at most
\(2/M_{\rm late}\). Any conditional label kernel supported on this
mask has expected price at most that value, for every entering history.
Separation therefore gives a fixed mixture of masks with baseline
marginal at most \(2/M_{\rm late}\) pointwise. Rows of the current
class do not enter one another's kernels, so their profiles can be
chosen independently. Use these same profiles for terminal avoidance.

\paragraph{Alarms and injective sampling at one class.}
After \(q\) classes have been completed, use the threshold
\[
 a_q=\exp\big(-n^\delta+q n^\delta/(2r)\big)
\]
for conditional reference probabilities of every still-future bad
\(F\). Terminal avoidance supplies \(a_0\). At the next class,
condition on its entering history and require both avoidance of the
current bads and that every future probability be at most \(a_{q+1}\).
Under its product reference law a future alarm costs at most
\(a_q/a_{q+1}=\exp(-n^\delta/(2r))\) by Markov; current bads
satisfy the incoming bound. The last threshold is still
\(a_r=e^{-n^\delta/2}\).

Let \(q_b(y\mid h)\) be the reference label marginal at history
\(h\), after integrating mask and sketch data. If
\(\sup_y\sum_bq_b(y\mid h)\le\theta_0\), apply the clock lemma
with parameter \(n_{\rm clock}=\lfloor e^{n^{\delta'}}\rfloor\),
where \(0<\delta'\ll\min(\delta,1)\), and \(B=5\).
The required comparisons are
\[
 \begin{gathered}
 \log M_{\rm late}=O(n)\ll n_{\rm clock},\qquad
 2e^{g_0n}/M_{\rm late}\ll n_{\rm clock}^{-A_*},\\
 n^{O(r)}\ll n_{\rm clock}^5,\qquad
 e^{-n^\delta/(2r)}\ll n_{\rm clock}^{-P_*}.
 \end{gathered}
\]
The second bound is the deterministic reference atom cap, including
fallbacks. The third covers both predicate scopes and row incidences,
by direct spatial locality with history fixed. The sampler assigns
distinct labels, avoids the current bads and future alarms, and has
joint upper comparison at factor at most two for the localized outputs
needed below. Side data are part of those outputs. If the column
condition fails, stop the run instead.

\paragraph{Column loads on reached histories.}
Fix a class and label \(y\). Its class size is
\(L_{\rm late}=2^{n-1}/H_*\). Put
\[
 Z_b=M_{\rm late}q_b(y\mid h),\qquad
 A_y=L_{\rm late}^{-1}\sum_b Z_b.
\]
Take the \(n\)-th moment with the indicator that the run has reached
this prehistory. A row weight has cap \(2e^{g_0n}\). Declare positions
near if their full baseline consultation scopes overlap, including
reference ancestors, deterministic initial resampling horizons and
slot variables. Scope and incidence counts give near fraction
\(f\le2^{-n+o(n)}\), so
\(nf\,2e^{g_0n}=o(1)\).

For at most \(n\) separated positions, integrate previous class
outputs backwards. At each step the nonnegative integrand is the
remaining product of queried conditional row masses. Retain the event
that the run reached that earlier stage with its incoming probability
and column bounds until its actual kernel has been compared with the
reference product. This costs one factor two for the entire queried
set. After the comparison, remove that stage's load check and all
later reach restrictions; retain only the reach condition needed to
compare the next earlier kernel. Reference kernels consult direct
local data, so the queried sets per stage remain of size \(n^{O(r)}\).

After at most \(r\) such comparisons, apply terminal comparison on
the required initial regions, then replace their permutation slots by
iid slots, at \(1+o(1)\) cost. The resulting law is exactly the
baseline in which profiles were balanced, including its finite
resampling. Separated consultation scopes give independent computations
in this baseline, each of normalized mean at most two. Thus scattered
moments give
\[
 \mathbb E[\mathbf1_{\rm reach}A_y^n]\le2^r K^n.
\]
Since the actual column sum is \((L_{\rm late}/M_{\rm late})A_y\),
\[
 \Pr\left(\mathrm{reach},\ \sum_bq_b(y\mid h)>\theta_0\right)
 \le2^r\left(\frac{K L_{\rm late}}{\theta_0M_{\rm late}}\right)^n.
\]
Here \(L_{\rm late}/M_{\rm late}\to0\). Union over the
\(\exp(O(n))\) label/class choices gives probability \(o(1)\)
that any reached stage stops. On a complete run all gated bads have
been avoided; induction from initial validity then supplies every
true-hit and deletion conclusion.
\end{proof}

\subsection{Matching the low-mode lists}
\label{sec:matching-the-low-mode-lists}

On a full run let \(q_v\) be the normalized final filtered law at
even role \(v\). Successive restrictions preserve its original cleaned
corner support and assigned palette. Draw an ordered pair from
\(q_v\otimes q_v\)
conditioned on nonconflict. Given the completed run, use independent
pair draws at different rows. The conditional pair law at one row is
\[
 \frac{q_v(x)q_v(z)\mathbf1_{\{|K_{\pi_i}(x,z)|\le\xi\}}}{Z_v},
 \qquad
 1-Z_v\le\Delta_i\|q_v\|_\infty
 \le8K^2C_n^{-1}e^{-198s_i}=o(1).
\]
Thus \(Z_v\ge1/2\) eventually. Each pair has distinct endpoints in
the row's assigned palette, and each endpoint hits every actual odd
neighbor.

Fix one patch and palette, writing \(M=M_i/\chi_i\). It has at
most \(2M\) labels. The allocated patch has
\(\Theta(2^nM_i/N)\) even roles, divided equally among its
\(\chi_i\) role colors by Lemma~\ref{lem:low-mode-palettes}.
Thus the even-role count in this palette satisfies
\(L\le KM/C_n\), and the low-mode size bounds give
\(L=2^{n-o(n)}\). Fix \(t'\le n\) distinct
roles in it. Let \(C(v)\) be the own cell of \(v\) together with
the cells supplying its external early labels. Join two tested roles
in the geometric overlap graph if their sets \(C(v)\) intersect or
they have a common late neighbor. Each role has at most
\(n^{A_c+O(1)}\) partners. Let \(j\) be this graph's rank, the
number of vertices minus the number of components, and \(m\le2j\)
the number of nonisolates. Enlarged resampling horizons need not be
disjoint: a joint comparison will remove that dependence before we
use disjointness of the direct cell sets.

\begin{proposition}[Joint endpoint kernels]\label{prop:low-endpoint-kernels}
For every fixed tuple of positions just described there are deterministic
nonnegative symmetric pair weights \(H_v\), chosen before endpoint
values, such that every assignment of two endpoints per row satisfies
\begin{equation}
 \Pr(\text{specified pairs, full run})\le e^{O(r)} K^{t'}
       \exp(.01 n m+C' t'j)\prod_v H_v(x_v,z_v),                  \label{eq:source-28}
\end{equation}
where \(C'\) is fixed and the nonnegative symmetric kernels on palette labels (allowing different kernels according to the positions) satisfy
\[
 \sup_x\sum_z H_v(x,z)\le K/M,\qquad \max H_v\le M^{-2}e^{.01n}.
\]
In particular we can use \(H_v=M^{-2}\) on nonisolates.
\end{proposition}

We specify the initial-state test in
Figure~\ref{fig:endpoint-comparisons}. Let \(J_v(S)\) assert initial
validity at \(v\): a valid internal prior with its cap and
single-corner support, \(S_v\) false, and palette mass at least
\(1/(4\chi_i)\). Let
\[
 A_v(x,z)=\mathbf1_{\{x,z\text{ in the assigned palette and }\mathcal X_i\}}
             \mathbf1_{\{|K_{\pi_i}(x,z)|\le\xi\}}.
\]
Define \(\Phi(S)=0\) if any \(J_v(S)\) fails; otherwise set
\[
 \Phi(S)=\chi_i^{2t'}\prod_v
 \left[A_v(x_v,z_v)
   \mathbf1_{\{x_v,z_v\in\mathcal D_v^0(S)\}}
   U_v^0(x_v;S)U_v^0(z_v;S)\right].
\]
The support indicator is redundant when the list product is positive,
but records precisely the support retained in the comparison.
This test reads only \(C_*:=\bigcup_v C(v)\), with at most
\(n(n+1)\) cells. Let \(\mathcal G\) require individual pool
typicality and \eqref{eq:source-23} throughout their deterministic
resampling horizons. Let \(\mathcal T\) retain only individual
pool typicality. These scopes fit the terminal and resampling
comparison propositions, even if their horizons overlap.

\begin{figure}[tbp]
\centering
\begin{tikzpicture}[
  box/.style={draw=black!65,rounded corners=2pt,fill=black!2,
    align=center,font=\small,inner sep=6pt,text width=12.1cm,
    minimum height=1.05cm},
  arrow/.style={-{Stealth[length=2mm]},line width=.65pt,draw=black!75},
  note/.style={font=\footnotesize,align=left,text width=5.6cm}
]
\node[box] (actual) at (0,0)
  {Specified endpoint pairs on a completed run\\
   Retain the events that each stage is reached and each pair test holds};
\node[box] (late) at (0,-2.3)
  {Independent reference outputs within each late class,\\
   conditional on its entering history};
\node[box] (lists) at (0,-4.6)
  {Initial assignment conditioned to avoid terminal failures\\
   Evaluate \(\mathbf1_{\mathcal G}\Phi(S_{\rm fin})\)};
\node[box] (initial) at (0,-6.9)
  {Initial assignment before conditioning on terminal avoidance\\
   Use permutation pools and finite resampling; retain \(\mathcal G\)};
\node[box] (fresh) at (0,-9.2)
  {Independent fresh cell states given the same fixed pools\\
   Use the calibrated cell strategies; retain \(\mathcal G\)};
\node[box] (iid) at (0,-11.5)
  {Fresh cell states with independent uniform pool slots\\
   Retain individual pool typicality \(\mathcal T\)};

\draw[arrow] (actual) -- (late)
  node[midway,right=3pt,note]
  {Bound the late outputs one class at a time, in reverse order\\
   Proposition~\ref{prop:late-assignment}};
\draw[arrow] (late) -- (lists)
  node[midway,right=3pt,note]
  {Integrate pair-hit factors with their side-data gates\\
   Bound~\eqref{eq:source-27}};
\draw[arrow] (lists) -- (initial)
  node[midway,right=3pt,note]
  {Conditioned expectation is at most\\
   \(1+o(1)\) times its preconditioning value\\
   Proposition~\ref{prop:terminal-avoidance}};
\draw[arrow] (initial) -- (fresh)
  node[midway,right=3pt,note]
  {Bound the same nonnegative test by fresh sampling at fixed pools\\
   Proposition~\ref{prop:initial-resampling-comparison}};
\draw[arrow] (fresh) -- (iid)
  node[midway,right=3pt,note]
  {First weaken \(\mathcal G\) to \(\mathcal T\), then compare slots\\
   Lemma~\ref{lem:late-classes-cells}};
\end{tikzpicture}
\caption{Upper comparisons for the endpoint product in
Proposition~\ref{prop:low-endpoint-kernels}.
The first two operations alternate, one class at a time in reverse order.
The late comparisons are conditional at each class, with the entering
history indicator retained until that comparison; each pair side-data
gate remains until its hit factor is integrated.
The pool gate \(\mathcal G\) includes individual typicality and
\eqref{eq:source-23} throughout the fixed deterministic simulation
horizons. It is retained through the fixed-pool resampling comparison,
then weakened to individual typicality \(\mathcal T\) before the slot
comparison. Every arrow bounds the indicated nonnegative test; none
asserts equality of laws or independence after avoidance.
The diagram stops before the later within-cell label, bin, pool and
history comparisons for nonisolated rows.}
\label{fig:endpoint-comparisons}
\end{figure}



\begin{proof}\leavevmode\par
\noindent
Figure~\ref{fig:endpoint-comparisons} records the comparison order for the initial-list reduction below.

\paragraph{Late factors and their side-data gates.}
Write \(a_v\ge1/(4\chi_i)\) for the initial palette mass and
\(m_{v,s}\ge.5-3e_{v,s}\) for each subsequent true hit mass.
On a full run the final row is
\[
 q_v(x)=\frac{\mathbf1_{\rm palette}(x)U_v^0(x)}{a_v}
       \prod_{s:\,v\sim b_s}\frac{I_x(y_{b_s})}{m_{v,s}}.
\]
Pair normalization costs at most two, and the squared hit denominators
cost \(4\exp(O(e_{v,s}))\) per late incidence. Summable errors
therefore bound the conditional product of specified pair weights by
\(K^{t'}\Phi(S_{\rm fin})\) times
\(\prod_{v\sim b_s}4I_{x_v}(y_{b_s})I_{z_v}(y_{b_s})\).

At a late role shared by several tested rows, retain one pair hit by
a rule fixed from the position tuple and replace the other factors
by four. Two distinct even roles have at most two common neighbors;
inside geometric components this removes \(O(t'j)\) incidences.
For each retained pair keep the side-data predicate certifying
\eqref{eq:source-27}. It reads the entering history, mask and sketches,
not the current label.

Integrate backwards by class. Keep that class's reach and incoming
bounds until applying its conditional sampler comparison at factor
two. The side-data predicate remains while the corresponding reference
label is integrated, so each retained factor \(4I_xI_z\) has
integral at most \(\exp(O(l e_{v,s}))\). Only then remove its
predicate. The product of these costs is \(K^{t'}\), since each row
has at most one incidence in each class and
\(\max(1,h_i)\sum_s e_{v,s}=O(1)\). The stage comparisons contribute
\(2^r=e^{O(r)}\), and removed incidences contribute \(e^{O(t'j)}\).
What remains is exactly the nonnegative initial test \(\Phi\)
defined above, with none of the late side-data gates left in it.

\paragraph{Terminal and resampling comparisons.}
Retain \(\mathcal G\) while comparing terminal-conditioned initial
sampling to the unconditioned process. Conditional on these pools,
Proposition~\ref{prop:initial-resampling-comparison} applies jointly
to the nonnegative test \(\Phi\) on \(C_*\), replacing final states
by independent fresh states at factor \(1+o(1)\). Only after this
fixed-pool comparison weaken \(\mathcal G\) to \(\mathcal T\) and
replace consulted permutation slots by iid slots.

Write \(\mathbb E_{\rm term}\) for terminal-conditioned initial
sampling, \(\mathbb E_{\rm perm}\) for the original process with
permutation pools and finite resampling, and
\(\mathbb E_{\rm fresh}(\,\cdot\mid\mathcal P)\) for independent
fresh states on fixed pools. The zero convention in the definition
of \(\Phi\) makes it a nonnegative test on all configurations.
These comparisons give
\[
\begin{aligned}
\mathbb E_{\rm term}\!\left[\mathbf1_{\mathcal G}\Phi(S_{\rm fin})\right]
&\le (1+o(1))\,
 \mathbb E_{\rm perm}\!\left[\mathbf1_{\mathcal G}\Phi(S_{\rm fin})\right]\\
&\le (1+o(1))\,
 \mathbb E_{\rm perm}\!\left[\mathbf1_{\mathcal G}
       \mathbb E_{\rm fresh}(\Phi\mid\mathcal P)\right]\\
&\le (1+o(1))\,
 \mathbb E_{\rm iid}\!\left[\mathbf1_{\mathcal T}
       \mathbb E_{\rm fresh}(\Phi\mid\mathcal P)\right],
\end{aligned}
\]
where the last expectation uses independent uniform slots, and each \(1+o(1)\) may absorb the preceding initial-comparison factors. The fixed-pool comparison is applied before weakening \(\mathcal G\) to \(\mathcal T\). The endpoint probability with full-run indicator is bounded by \(\exp(O(r)+O(t'j))K^{t'}\) times the left side, including the cost of omitted shared late incidences in the backward integration above.

\paragraph{Isolated rows.}
After the three displayed comparisons, discard typicality conditions
on pools outside \(C_*\). The direct cells of an isolated row are
disjoint from every other row's cells, so their iid-pool fresh
expectation factors. An upper pair weight is
\[
 H_v(x,z)=\chi_i^2 A_v(x,z)\,
 \mathbb E_{\rm iid,fresh}\!\left[
 \mathbf1_{\mathcal T_v}\sigma_v(x)\sigma_v(z)
 \prod_{b\in B_v^e}
       \frac{I_x(y_b)I_z(y_b)}{D_{i(b)}(x)D_{i(b)}(z)}\right],
\]
where \(\mathcal T_v\) requires typicality of the direct pools,
and invalid internal-prior contributions are zero. We have removed
the initial mass, masked-support and palette-retention gates and the
redundant positive-support indicator from the row factor of \(\Phi\).
Every removal increases the integral.

The external observation cells are distinct and separate from the
own cell. Their fresh singleton comparisons give
\((1+O(n^{-3}))\pi_{i(b)}\) after integration. Then
\eqref{eq:source-24} and Lemma~\ref{lem:fresh-internal-prior} give
\[
 \sum_z H_v(x,z)
 \le K\chi_i\,\mathbb E[\mathbf1_{\rm typ}\sigma_v(x)]
 \le K\chi_i/M_i=K/M.
\]
For the entrywise bound, each prior atom is at most
\(M_i^{-1}e^{O(\log n)}\); crossing factors contribute
\(e^{O(\log n)}\), and the nonconflict cutoff and degrees bound
bulk pair-hit ratios by \(1+O(\xi+\log n/n)\). Hence
\(H_v\le M^{-2}\exp(O(\log n)+O(\xi n))\le M^{-2}e^{.01n}\)
by the prior choice of \(\xi\).

\paragraph{Reducing nonisolates to pair-hit queries.}
For each of the other \(m\) rows first use, on \(J_v\), the cap
\(\sigma_v\le c_i:=M_i^{-1}e^{O(\log n)}\). In particular
\[
 \mathbf1_{J_v}\mathbf1_{\{x,z\in\mathcal D_v^0\}}
                   \sigma_v(x)\sigma_v(z)\le c_i^2.
\]
After this inequality remove all remaining state-dependent
initial-validity, support, mass and palette-retention indicators for
that row. Keep only the deterministic endpoint predicate \(A_v\).
Bound crossing denominators on the envelope and discard their hits.
Bulk pair denominators cost \(4(1+O(\log n/n))\) each. The remaining
random factors are pair-hit indicators of external early bulk labels.
There are no residual tests of the removed own priors.

Index such a query by its even row, observed odd role, outer word,
observed slice, cell and group. At each observed slice choose one even
outer-word family to retain, by a rule depending only on the positions.
Discard the other hit incidences there. Two different outer words have
at most two common one-flip words, and a row has at most one observation
in a fixed slice. Charging to pairs of rows within the geometric
components removes \(O(t'j)\) hits. Let \(k\le mn\) be the number
of retained queries.

Queries in the same slice now arise from one even outer word. Their
internal words have the same palette color and are distinct, so their
distances exceed \(500\rho h_i\). In cluster mode two members of
one projected group differ in at most two internal bits, whereas the
fixed threshold ensures \(500\rho h_i>2\). Thus these groups are
distinct. Their raw primitive consultation radii are at most
\(10\rho h_i\), so those consultations are disjoint as well. When
\(h_i=0\) there is no sharing. This independence will be used only
after removing slice conditioning below.

\paragraph{Repeated bins and reverse summation.}
Apart from deterministic factors, the integral to bound is now
\[
 \mathbb E_{\rm iid,fresh}
       \left[\mathbf1_{\mathcal T_{\rm qry}}\prod_{a=1}^k f_a(Y_a)\right],
 \qquad f_a(y)=I_{x_{v(a)}}(y)I_{z_{v(a)}}(y),
\]
where \(\mathcal T_{\rm qry}\) requires typical pools for the
queried odd cells. Conditional on these pools and their successful
\(W\), let \(D_a\) be the bin chosen for group \(g_a\), and set
\[
 u_a^W(D)=\int f_a\,dU_{g_a}^W(D),\qquad
 J(D)=\{a:D_a=D_b\text{ for some }b<a\text{ in the same cell}\}.
\]
For fixed bin values discard \(f_a\le1\) whenever \(a\in J(D)\).
The remaining queries use at most one label per bin in each cell, so
the conditional label comparison costs \(e^{d_i^{-.01}k}\).
Next the group-bin comparison, applied to all \(k\) distinct group
queries, costs \(e^{d_i^{-.05}k}\). In direct mode bins are singleton
labels; its calibrated cell sampler gives the corresponding upper
comparison at \(e^{o(1)k}\), since \(k\le n^2\).

For a fixed repeat mask \(J\), the sum after these comparisons is
bounded by the appropriate calibration factor times
\[
 \sum_{D_1,\ldots,D_k}\mathbf1_{\{J(D)=J\}}
          \prod_{a=1}^k\widetilde q_{g_a}(D_a)
          \prod_{a\notin J}u_a^W(D_a).
\]
Enlarge the repeat indicator by retaining only distinctness of the
kept targets within each cell and the necessary earlier-bin equality
for each removed query. Sum removed bins in decreasing index order.
At the summation for \(a\in J\), all earlier candidate bins are
still fixed, even if they will be removed later, so
\[
 \sum_{D_a\in\{D_b:b<a\text{ in the same cell}\}}
                   \widetilde q_{g_a}(D_a)
 \le k\|\widetilde q_{g_a}\|_\infty\le n^{-A_c+3}.
\]
No kept factor \(u_a^W(D_a)\) reads a removed bin. This is why the
initial indicators were discarded before the two sampler comparisons.
The total repeat cost is \((n^{-A_c+3})^{|J|}\); kept-bin
distinctness still remains for the pool integral.

\paragraph{Removing pool restrictions and counting the losses.}
On typical pools the kept restricted bin probabilities satisfy
\[
 \widetilde q_g(D)\le
 (1+\varepsilon)(B_i/L_i^{\rm slot})q_g^W(D)
                     \mathbf1_{\{D\in\mathcal P\}},
 \qquad \varepsilon=10^{-5}.
\]
Here \(q_g^W\) is the raw bin law, uniform before trims in direct
mode; the bound includes both trim losses. Undo the \(W\)-cell
load conditioning, at negligible logarithmic cost per query, while
its typical-pool hypothesis is still retained. Then discard the
remaining pool tests. The slice-conditioned \(W\)'s are now
independent of the iid pools. For \(a\) distinct kept targets in
one cell, their simultaneous presence costs at most
\[
 (L_i^{\rm slot})_a/B_i^a\le(L_i^{\rm slot}/B_i)^a.
\]
This cancels the pool-restriction factors. Discard kept-target
distinctness and sum their raw bin kernels. Finally remove the
\eqref{eq:source-15} slice conditioning at cost at most another
\((1+\varepsilon)^k\). Raw primitive inputs are now independent
between slices and between the separated consultations within them.

Their pair-hit means use \(\pi_i^0\) in cluster mode, and are at
most \((1+10^{-4})/4\), by the degree bounds, the pair cutoff and
the total-variation estimate following \eqref{eq:source-16}.
Let \(\delta_{\rm cal}\) include the two per-query calibration
logarithms and the load-normalizer logarithm. With \(\xi\) already
fixed small, increase the fixed threshold \(Q_0\) so that
\[
 \delta_{\rm cal}+2\log(1+\varepsilon)
       +\log(1+10^{-4}+4n^{-A_c+3})\le .005
\]
eventually. This makes the fixed low-cluster errors small; it does
not assume they tend to zero with \(n\). Summing all repeat masks
therefore bounds the \(k\) pair hits by
\[
 e^{(\delta_{\rm cal}+2\log(1+\varepsilon))k}
       \big((1+10^{-4})/4+n^{-A_c+3}\big)^k
 \le4^{-k}e^{.005k}.
\]
For \(k=0\) the bound is one. Each retained pair hit contributes
its factor \(1/4\), cancelling a bulk denominator factor four.
Discarded hits leave \(e^{O(t'j)}\); the prior caps, crossing
factors and degree drift leave \(e^{O(m\log n)}\). Since
\(k\le mn\), the total is at most
\(M_i^{-2m}\exp(.01nm+O(t'j))\). The palette factor
\(\chi_i^{2m}\) turns this into \(M^{-2m}\) with the asserted
error. Together with the isolate weights and the preceding late and
initial comparisons, this proves \eqref{eq:source-28}.
\end{proof}

\begin{lemma}[Two-endpoint Hall estimate]\label{lem:low-two-endpoint-hall}
Assume that the full run has probability \(1-o(1)\), as supplied by
Proposition~\ref{prop:late-assignment}. The joint endpoint bounds of
Proposition~\ref{prop:low-endpoint-kernels}, with the stated palette sizes
and geometric partner counts, then imply simultaneous Hall matchings in
every low-mode palette with positive probability of full success.
\end{lemma}

\begin{proof}\leavevmode\par
\paragraph{Geometric averaging.}
Choose a small constant \(\eta>0\) with
\(.02+C'\eta<(\log2)/2\), and put \(t_0=\lfloor\eta n\rfloor\).
For a uniformly chosen distinct \(p\)-set of positions, \(p\le t_0\),
a geometric overlap graph of rank \(j\) contains a forest of \(j\)
necessary partner relations. Specifying its edges costs at most
\(p^{2j}\); exposing along that forest uses the partner fraction
\(2^{-n+o(n)}\) once per edge. Since \(m\le2j\), the correction in
\eqref{eq:source-28} is at most
\(e^{(.02+C'\eta)nj}\). Its average is bounded by a convergent
geometric sum with ratio \(e^{-(\log2)n/2+o(n)}\).

\paragraph{The endpoint multigraph.}
This graph is different from the geometric overlap graph: its vertices
are host labels, and each even row is an edge joining its two chosen
endpoints. Hall fails precisely when some set of row edges has fewer
incident endpoint vertices than edges. Such a set has a connected
violating component. The actual distinct-endpoint condition is retained
in the following sums, even for nonisolates whose convenient upper
weight \(H_v=M^{-2}\) is positive on the diagonal.

Connectivity of \(p\) distinct row edges is witnessed by a tree on
those rows, with one endpoint identification along each tree link.
There are at most \(4^{p-1}p^{p-1}\) such diagrams. These identifications
leave \(p+1\) abstract endpoint variables, and the row edges form a
tree on them. Summing a product of the row weights by successively
integrating leaves costs at most \(2M(K/M)^p\), using the row-sum
bound and at most \(2M\) choices for the last variable.

A connected Hall-violating \(p\)-set, \(p\ge3\), uses fewer than
\(p\) endpoint vertices. Thus at least two further mergers of distinct
current endpoint classes occur, with \(O(p^4)\) choices. Keep just
these two mergers and discard any additional identifications for an
upper bound. The resulting quotient has \(p-1\) vertices and \(p\)
row edges. On the retained distinct-endpoint support it has no loops;
parallel edges are allowed. A spanning tree uses \(p-2\) edges,
leaving exactly two edges to which the maximum-weight bound applies.
The corresponding sum is at most
\[
 \underbrace{2M}_{\text{root choices}}\,
 \underbrace{(K/M)^{p-2}}_{\text{tree edges}}\,
 \underbrace{(M^{-2}e^{.01n})^2}_{\text{two remaining edges}}
 \le K^pM^{-p-1}e^{.02n}.
\]
The mergers here are independent identifications of equivalence classes,
not probabilistically independent events. Different rows may use
different kernels; the same two uniform bounds suffice.

\paragraph{Summing obstructions.}
Equation~\eqref{eq:source-28}, geometric averaging and
\(\binom Lp\le L^p/p!\) bound the expected number, with full-run
indicator, of connected \(t_0\)-row sets by
\(e^{O(r)}M(K/C_n)^{t_0}\). For smaller connected \(p\)-row sets
with fewer than \(p\) endpoints the bound is
\[
 e^{O(r)}p^4M^{-1}e^{.02n}(K/C_n)^p.
\]
Sizes one and two cannot violate Hall because each row has distinct
endpoints. For the large-set estimate,
\[
 \log\big(M(K/C_n)^{t_0}\big)
 =n\log2+o(n)+\eta n\log(K/C_n)+O(\log n),
\]
which is eventually negative at a linear rate, however slowly
\(C_n\to\infty\). For the small-set sum,
\(M^{-1}e^{.02n}=e^{-(\log2-.02)n+o(n)}\), and
\(\sum_{p\ge3}p^4(K/C_n)^p\) is bounded. Both estimates still tend
to zero after summing over the \(e^{o(n)}\) patches and palettes.

Every connected component with at least \(t_0\) row edges contains
a connected \(t_0\)-row set. Every Hall failure contains a connected
violating component. Consequently the probability of a full run with
no Hall obstruction is at least its probability \(1-o(1)\) minus
the preceding expected obstruction counts, and is positive. Hall gives
distinct representatives simultaneously in all palettes. Palettes
and patches are disjoint, so these choices combine. The early odd
labels are already injective, and the disjoint injective late pools
replace their dummy labels. Every required edge is present, yielding
the monochromatic cube and excluding the low modes.
\end{proof}

\paragraph{Conclusion of Theorem~\ref{thm:main}.}
An unbounded ratio \(R(Q_n)/2^n\) would supply the sequence of
Lemma~\ref{lem:counterexample-sequence}. The earlier exclusions reduce
it to deep discrepancy in both orientations. Patch extraction then
leaves high or low modes. Section~\ref{sec:high-modes-of-the-tiling}
excludes the high modes, and the preceding matching excludes the low
modes. No such sequence exists. The ratio therefore has finite
supremum, giving an absolute \(C\) with \(R(Q_n)\le C2^n\) for all
\(n\). This existence argument does not provide a numerical value
for \(C\).
